\documentclass[10pt,a4paper]{amsart}
\usepackage[margin=19mm]{geometry}
\usepackage{amsmath,amssymb,mathtools}
\usepackage[scr=rsfs,cal=euler]{mathalfa}
\usepackage{xfrac}
\usepackage[unicode,hidelinks,bookmarks=false]{hyperref}
\newcommand{\ie}{i.e.\ }
\newcommand{\cf}{cf.\ }
\newcommand{\resp}{respectively\ }
\newcommand{\Subsection}{\subsection}
\providecommand{\nobreakdash}{\nobreak\hbox{-}}
\setlength{\emergencystretch}{2em}
\multlinegap0pt
\usepackage{mathtools}
\usepackage{amssymb}
\usepackage{xcolor}
\usepackage{float}
\usepackage{tikz}
\usepackage[shortlabels]{enumitem}
\newtheorem{theorem}{Theorem}
\title{Barta Theorem for the $p$-Laplacian \\and Geometric Applications}
\author{Paulo Henryque C. Silva}
\date{Source: arXiv:2603.08350v2 (CC BY 4.0)}
\begin{document}
\maketitle

\noindent\emph{Source and licence.} Barta Theorem for the $p$-Laplacian \\and Geometric Applications. Version arXiv:2603.08350v2 (CC BY 4.0), distributed by arXiv under the Creative Commons Attribution 4.0 International licence, http://creativecommons.org/licenses/by/4.0/, as recorded in the arXiv licence metadata for this exact version and shown on the abstract page https://arxiv.org/abs/2603.08350v2. Full licence notice: \texttt{licenses/arxiv-2603.08350-CC-BY-4.0.txt}.\par\smallskip

\noindent\textit{Editorial scope.} Counted unit: the article's theorem thm:pminsub\_ball (source0.tex lines 281--295). The article's complete original proof of it is delivered with it (source0.tex lines 1358--1388, one \texttt{proof} environment of 31 lines, which the article places away from the statement). No mathematics is written, repaired or re-derived: the statement and the proof are the article's own lines, in the article's own notation, and no retained line is substituted or re-flowed.  Retained uncounted as the source-local context the proof needs: lines 127--131; lines 1353--1355.  Nothing was omitted from any retained span, so the package carries no omission record.  No retained line contains a citation, so the wrapper carries no bibliography and no bibliography entry.  No retained line contains a figure environment, an image-inclusion command or drawing code, so no image is delivered and the package declares no asset dependency.  Editorial formatting only, in addition: an AMS \texttt{amsart} preamble that restates the article's own declarations and macros used by the retained lines, namely the environments \texttt{theorem} on separate counters, together with the standard packages those lines need.  The retained environments print in this document as Theorem 1 in order of appearance, whereas the article prints the same items with the numbering of its own full document.  The complete native source of the article as distributed in the arXiv e-print for this exact version is bundled unchanged as \texttt{sources/195996/source0.tex}.
\begin{equation}\label{eq:Rayleigh quotient}
    \lambda_{1,p}^{\mathcal{D}}(\Omega)
    = \inf \left\{ \frac{\displaystyle \int_{\Omega} |\nabla \varphi |^p d\mu} {\displaystyle \int_{\Omega} |\varphi|^p d\mu}; \, \varphi \in \mathrm{W}^{1,p}_0(\Omega), \, \varphi \not\equiv 0
    \right\}.
\end{equation} Clearly,  \eqref{eq:Rayleigh quotient} is valid for  arbitrary open subsets $\Omega \subseteq M$ and defines what is called the \emph{$p$-fundamental tone} $\lambda_{1,p}(\Omega)$ of  $\Omega $, although $\Omega$ may not have Dirichlet eigenvalues.

\begin{equation}\label{lowerradiostone}
r \geq \left( k^{p-2} \cdot \frac{\lambda_{1,p}^{\mathcal D}(B_1^m)}{\lambda_{1,p}(\Omega)}\right)^{1/p},
\end{equation}

\begin{theorem}\label{thm:pminsub_ball}
Let $(N^n,g)$ be an $n$-dimensional Riemannian manifold and fix $p_0\in N$. Assume that the radial sectional curvatures of $N$ with respect to $p_0$ satisfy $K_N(\partial_t,v) \leq c$ for every $x\in B_N(p_0,r)\setminus \mathrm{Cut}(p_0)$ and every unit vector 
$v\perp \partial_t$. Suppose moreover that $\mathcal{H}^{n-1}\!\big(\mathrm{Cut}(p_0)\cap B_N(p_0,r)\big)=0$. Let $\varphi:M^m\hookrightarrow N^n$ be a minimal immersion and let 
$\Omega\subset \varphi^{-1}(B_N(p_0,r))$ be a connected component. 
Assume that the extrinsic distance function $t(x)=\operatorname{dist}_N(\varphi(x),p_0)$ satisfies $|\nabla^M t|\ge k>0$ in $\Omega$. Let $B(r)\subset \mathbb{M}^m(c)$ denote the geodesic ball of radius $r$ in the 
$m$-dimensional simply connected space form of constant curvature $c$. 
Then for every $2\le p<\infty$ and $r\le r_\ast(c)$, it holds that
\[
\lambda_{1,p}(\Omega)
\ge
k^{p-2} \cdot \lambda_{1,p}^{\mathcal D}(B(r)).
\]
If equality holds for some bounded $\Omega$, then necessarily 
$k=1$ and $\Omega$ is isometric to $B(r)$.
\end{theorem}

\begin{proof}
Let $B_r^{m}$ denote the Euclidean ball of radius $r$ in $\mathbb{R}^{m}$. To prove the desired result, it suffices to show that
\begin{equation}\label{scaling}
    \lambda_{1,p}^{\mathcal{D}}(B^{m}_r) = \frac{1}{r^p}\lambda_{1,p}^{\mathcal D}(B_1^m).
\end{equation}
Indeed, if \eqref{scaling} holds, then Theorem \ref{thm:pminsub_ball}
yields the following lower bound for the $p$-fundamental tone
\begin{equation}
    \lambda_{1,p}(\Omega) \geq k^{p-2} \cdot \lambda_{1,p}^{\mathcal{D}}(B^{m}_r) = k^{p-2} \cdot \frac{1}{r^p}\lambda_{1,p}^{\mathcal D}(B_1^m).
\end{equation}
Thus establishing the desired inequality in \eqref{lowerradiostone}.
It remains to prove \eqref{scaling}. To this end, we use the variational
characterization given in \eqref{eq:Rayleigh quotient}. Consequently,
\begin{equation}
      \lambda_{1,p}^{\mathcal{D}}(B^{m}_r)
    = \inf \left\{ \frac{\displaystyle \int_{B^{m}_r} |\nabla u |^p dx} {\displaystyle \int_{B^{m}_r} |u|^p dx}; \, u \in \mathrm{W}^{1,p}_0(\Omega), \, \varphi \not\equiv 0
    \right\}.
\end{equation}
Now consider the change of variables $x=ry$ and define $v(y)=u(ry)$. Then
$\nabla_y v(y)= r\, \nabla_x u(ry)$, and, since $dx = r^{m}dy$, it follows that 
\begin{equation}\label{eqcor}
     \frac{\displaystyle \int_{B^{m}_r} |\nabla u |^p dx} {\displaystyle \int_{B^{m}_r} |u|^p dx} =  r^{-p} \cdot\frac{\displaystyle \int_{B^{m}_1} |\nabla v |^p dy} {\displaystyle \int_{B^{m}_1} |v|^p dy}.
\end{equation}
Finally, note that 
\[
T\colon W^{1,p}_0(B_r^{m}) \to W^{1,p}_0(B_1^{m}), 
\qquad (Tu)(y) = u(ry),
\]
is a bijective mapping. Hence, taking the infimum over 
$W^{1,p}_0(B_r^{m})$ in \eqref{eqcor} gives the equality in \eqref{scaling}.
\end{proof}

\end{document}
